% ============================================================================
% 07-integrality.tex — the arithmetic separation. Claims #7-#10 of PLAN.md §3.
% Attribution discipline: integrality of the s7 partner is O'Brien 2016 Thm 6.2
% (CITED), used in the formalization as the single registered literature axiom;
% the recurrence-matching step is kernel-checked. Nothing here re-proves it.
% ============================================================================
\section{The arithmetic separation of the three candidates}
\label{sec:integrality}

Theorem~\ref{thm:sym2-template} is blind to the choice of parameters. This
section exhibits the invariant that is not: the integrality of the holomorphic
solution of the order-2 partner. Among Cooper's three sporadic sequences, only
$s_7$ has an integral partner, and we make the mechanism behind that
integrality precise.

\subsection{The partner recurrence}

Extracting the coefficient of $z^{M}$ from $\Ltwo f = 0$ for
$f = \sum_{n\geq0}a_nz^n$ with $a_0 = 1$, using the partner
coefficients~\eqref{eq:generic-partner}, gives $a_1 = \tfrac{b}{2}$ and, for
$k\geq0$,
\begin{equation}\label{eq:partner-recurrence}
  (k+2)^2 a_{k+2}
  \;=\; a_{k+1}\Bigl(2a(k+1)^2 + a(k+1) + \tfrac{b}{2}\Bigr)
      - a_k\Bigl(ck^2 + ck + \tfrac{c+d}{4}\Bigr).
\end{equation}
The leading coefficient $(k+2)^2$ is the crux: the sequence is
$\Q$-valued by construction, and integrality --- when it holds --- is an
Ap\'ery-style theorem, not a consequence of the shape of the recursion.

There is, however, a floor. Let $s\colon\mathbb{N}\to\Z$ be \emph{any} integer
sequence with $s(0)=1$ and let $g$ be the coefficient sequence of the formal
square root of $\sum_n s(n)z^n$, defined by $g_0=1$ and
$g_n = \bigl(s(n) - \sum_{k=1}^{n-1}g_kg_{n-k}\bigr)/2$.

\begin{proposition}[Dyadic baseline; status (K)]\label{prop:dyadic}
For every integer sequence $s$ with $s(0)=1$, every $g_n$ lies in
$\Z[\tfrac12]$.
\end{proposition}

\begin{proof}
Induction on $n$: $\Z[\tfrac12]$ is a ring containing $\Z$, so the sum of
products of earlier terms lies in it, and it is closed under division by $2$.
Formalized as \artifact{sqrtSeq\_dyadic} in
\artifact{Agora/Sequences/FormalSqrt.lean}; the accompanying control
\artifact{dyadic\_baseline\_control} checks that $\tfrac13\notin\Z[\tfrac12]$,
so the statement is not vacuous.
\end{proof}

Proposition~\ref{prop:dyadic} concerns the formal square root $g$; the object
of interest is the recurrence-defined partner $a$
of~\eqref{eq:partner-recurrence}. That the two coincide is the
\emph{sequence-level} form of Theorem~\ref{thm:sym2-template}, and it is not a
formality: the operator identity $\Lthree = P_2\cdot\Symtwo(\Ltwo)$ says that
$\Symtwo(\Ltwo)$ annihilates $f^2$, and passing from there to a statement
about coefficient sequences is exactly the step that a proof assistant does
not grant for free.

\begin{theorem}[Sequence-level symmetric square; status (K)]\label{thm:bridge}
Let $(a,b,c,d)\in\Z^4$ and let $s\colon\mathbb N\to\Z$ satisfy the Cooper
recurrence
\begin{equation}\label{eq:cooper-recurrence}
  (n+1)^3 s(n+1) = (2n+1)\bigl(an^2+an+b\bigr)s(n) - n\bigl(cn^2+d\bigr)s(n-1)
  \qquad(n\geq1),
\end{equation}
the coefficient form of $\Lthree(a,b,c,d)\sum s(n)z^n=0$, with $s(0)=1$,
$s(1)=b$. Let $(a_n)$ be
the partner sequence defined by~\eqref{eq:partner-recurrence} with $a_0=1$,
$a_1=\tfrac b2$, and $(g_n)$ the formal square root of $\sum s(n)z^n$. Then
$a_n = g_n$ for every $n$. Equivalently
\[
  \sum_{i=0}^{n} a_i\,a_{n-i} \;=\; s(n)\qquad(n\geq 0).
\]
\end{theorem}

\begin{proof}
Write $y_n=\sum_{i=0}^n a_ia_{n-i}$. Three steps.

\emph{(i) Symmetrization.} For any kernel $F$ and any sequence $a$,
$\sum_{i=0}^{m}F(i,m-i)\,a_ia_{m-i}
  = \sum_{i=0}^{m}\tfrac12\bigl(F(i,m-i)+F(m-i,i)\bigr)a_ia_{m-i}$,
by reflecting the summation index.

\emph{(ii) The convolution satisfies the order-3 recurrence.} Multiply
\eqref{eq:partner-recurrence} at index $t=n-i$ by $(6i+2t+4)\,a_i$ and sum
over $0\le i\le n$. The three resulting sums have kernels
$(6j+2k)k^2$, \ $(6j+2k+2)(2ak^2+ak+\tfrac b2)$ and
$(6j+2k+4)(ck^2+ck+\tfrac{c+d}4)$ on the antidiagonals $j+k=n+2,\,n+1,\,n$
respectively, and their symmetrizations are \emph{constant} on each
antidiagonal, equal to
\[
  (n+2)^3,\qquad (2n+3)\bigl(a(n+1)^2+a(n+1)+b\bigr),\qquad
  (n+1)\bigl(c(n+1)^2+d\bigr).
\]
These are precisely the coefficients of the Cooper recurrence. The boundary
terms $i=n+1,n+2$ cancel using $a_1=\tfrac b2a_0$. Hence $(y_n)$ satisfies the
Cooper recurrence.

\emph{(iii) Uniqueness.} $y_0=1=s(0)$ and $y_1=2a_0a_1=b=s(1)$, and the leading
coefficient $(n+2)^3$ never vanishes, so $y=s$. Finally $\sum a_nz^n$ and
$\sum g_nz^n$ are two power series over $\Q$ with constant term $1$ and equal
squares; since $\Q[[z]]$ is a domain and their sum is a unit, they are equal.

Formalized as \artifact{partner\_eq\_sqrt} in
\artifact{Agora/Sequences/SqrtBridgeGeneric.lean}, with step~(ii) as
\artifact{conv\_cooper\_of\_rec\_generic} over arbitrary
$(A,B,C,D)\in\Q^4$; the weight $6i+2t+4$ is the entire creative content, and
the kernel verifies the rest by \artifact{ring}. The instances are
\artifact{partner\_eq\_sqrt\_s7} and \artifact{partner\_eq\_sqrt\_s10}, which
discharge the recurrence hypothesis by the WZ certificates of
\S\ref{sec:formalization}. No instance is stated for $s_{18}$: the development
contains no closed form for $s_{18}$ to instantiate at, and we do not assume
one.
\end{proof}

\begin{corollary}[status (K)]\label{cor:partner-dyadic}
Under the hypotheses of Theorem~\ref{thm:bridge}, every partner coefficient
$a_n$ lies in $\Z[\tfrac12]$. In particular this holds for the $s_7$ and
$s_{10}$ partners at every index.
\end{corollary}

So powers of $2$ in the denominators are the \emph{forced} baseline for the
partner of any Cooper-template instance, now as a theorem rather than as an
observation: $s_{10}$ exhibits it (Theorem~\ref{thm:non-integral} shows the
prime $2$ genuinely occurs), and it is $s_7$ that is anomalous. For $s_7$ the
corollary excludes every odd prime from the denominators unconditionally; it
does \emph{not} give integrality, which remains the content of
Corollary~\ref{cor:s7-integral} and rests on the citation discussed there.

The 2-adic half can be isolated exactly.

\begin{proposition}[status (K)]\label{prop:mod4}
Let $s\colon\mathbb N\to\Z$ with $s(0)=1$ and $4\mid s(n)$ for all $n\geq1$. Then
every coefficient $g_n$, $n\geq1$, of the formal square root is an even integer.
Consequently, if $4\mid s_7(n)$ for all $n\geq1$, the $s_7$ partner is integral,
with no appeal to Theorem~\ref{thm:obrien}.
\end{proposition}

\begin{proof}
Strong induction: if $g_k\in2\Z$ for $1\le k<n$ then
$\sum_{k=1}^{n-1}g_kg_{n-k}\in4\Z$, so
$g_n=\tfrac12\bigl(s(n)-\sum g_kg_{n-k}\bigr)\in2\Z$. The second statement is
Theorem~\ref{thm:bridge}. Formalized as \artifact{sqrtSeq\_even\_of\_four\_dvd}
and \artifact{s7\_partner\_integral\_of\_congruence} in
\artifact{Agora/Sequences/SqrtIntegrality.lean}.
\end{proof}

\begin{remark}[PASS(200)]\label{rem:pass200}
The hypothesis $4\mid s_7(n)$ holds for $1\leq n\leq200$ in exact integer
arithmetic, the minimum 2-adic valuation over that range being exactly $2$; it
is kernel-checked for $n\leq6$. This is evidence, not proof, and
Corollary~\ref{cor:s7-integral} does not rely on it. If proved, it would replace
the development's only literature axiom by an elementary congruence on the
binomial sum defining $s_7$. The criterion is sufficient but not necessary
($(1+z)^2$ has an integral square root), so it is not a characterization; it
does fail for $s_{10}$ already at $n=1$ (\artifact{s10\_fails\_criterion}).
\end{remark}

\begin{remark}[A cautionary note on ``blocked'' goals]\label{rem:blocked}
For $s_7$, Theorem~\ref{thm:bridge} stood for some weeks as the development's
last open goal, with four recorded failed strategies and an internal ruling
that it was blocked on the absence, in the pinned library, of a power-series
square root and of a holonomic-sequence API. Both absences are real and neither
was relevant: the proof above uses only reindexing of finite sums and the fact
that $\Q[[z]]$ is a domain. We record this because the failure mode is general.
A ``blocked on the library'' verdict is a statement about a proof route; four
failed strategies are evidence about four strategies.
\end{remark}

\subsection{Non-integrality for \texorpdfstring{$s_{10}$}{s10} and
\texorpdfstring{$s_{18}$}{s18}}

\begin{theorem}[status (K)]\label{thm:non-integral}
The holomorphic solution of the $s_{10}$ partner has coefficients beginning
$1,\;1,\;\tfrac{17}{2},\;\tfrac{147}{2}$, and that of the $s_{18}$ partner
begins $1,\;3,\;\tfrac{45}{2},\;\tfrac{429}{2}$. In particular neither
sequence is integral.
\end{theorem}

\begin{proof}
Evaluate~\eqref{eq:partner-recurrence} at the respective parameters; the terms
listed are exact rationals. Integrality of the $n=2$ term would force
$17 = 2m$ resp. $45 = 2m$ for an integer $m$, which is impossible. The
formalization states the refutation in this arithmetic form rather than as
denominator bookkeeping: \artifact{s10\_partner\_two\_not\_integral},
\artifact{s18\_partner\_two\_not\_integral}, and the conclusions
\artifact{s10\_partner\_not\_integral}, \artifact{s18\_partner\_not\_integral}
in \artifact{Agora/Sequences/PartnerIntegrality.lean}.
\end{proof}

Note the asymmetry of difficulty: refuting integrality needs one witness,
establishing it needs all of them. The formalization keeps that asymmetry
visible in the type of its summary statement
(\artifact{cooper\_partner\_separation}), which asserts only the two negative
halves.

\subsection{Integrality for \texorpdfstring{$s_7$}{s7}}

For $(a,b,c,d) = (13,4,-27,3)$, recurrence~\eqref{eq:partner-recurrence}
becomes, after clearing,
\begin{equation}\label{eq:s7-partner-recurrence}
  (k+2)^2a_{k+2}
  = \bigl(26(k+1)^2 + 13(k+1) + 2\bigr)a_{k+1}
    + 3(3k+1)(3k+2)\,a_k,
  \qquad a_0=1,\ a_1=2 .
\end{equation}
Its first terms are
\[
  1,\;2,\;22,\;336,\;6006,\;117348,\;2428272,\;52303680,\;\dots
\]
which is OEIS A279619~\cite{OEIS}. Integrality of this sequence is a theorem
of the literature, and we state it as such.

\begin{theorem}[cited, O'Brien 2016]\label{thm:obrien}
\cite[Thm.~6.2, p.~47]{OBrien2016}. Any rational sequence satisfying
recurrence~\eqref{eq:s7-partner-recurrence} with the stated initial values is
integer-valued.
\end{theorem}

O'Brien's proof is an explicit coefficient-matching induction using the
$q$-expansions of two specific weight-1 level-7 modular objects; it is not a
formal consequence of any general integrality principle, and \S\ref{ssec:mechanism}
explains why. We do not reprove it.

\begin{corollary}[status (K), modulo Theorem~\ref{thm:obrien}]
\label{cor:s7-integral}
The holomorphic solution of the $s_7$ partner is integral, and its coefficient
sequence is A279619.
\end{corollary}

\begin{proof}
That the sequence defined by the generic
recurrence~\eqref{eq:partner-recurrence} at the $s_7$ parameters satisfies
exactly~\eqref{eq:s7-partner-recurrence}, with the same initial values, is a
mechanical identity; it is kernel-checked as
\artifact{partnerSeq\_s7\_recurrence}. Theorem~\ref{thm:obrien} then applies.
In the formalization the cited theorem is the registered axiom
\artifact{obrien2016\_theorem6\_2}, and the conclusion is
\artifact{s7\_partner\_integral}; see \S\ref{sec:formalization} for the axiom's
statement and its disclosure.
\end{proof}

\begin{remark}[PASS(7)]\label{rem:pass7}
The formalization independently retains a kernel-checked finite verification,
\artifact{s7\_partner\_integral\_pass7}: the coefficients are integers for
$n\leq7$, each obtained by evaluating the recurrence rather than by table
lookup, with the eight witnesses displayed above. This was the honest position
before the citation was traced to its correct source, and it is kept on the
record. It is evidence, not proof, of the general statement, and is reported
as \emph{PASS(7)} rather than as a theorem about all $n$
(\S\ref{sec:formalization}).
\end{remark}

\subsection{The mechanism: a normalized integral uniformizer}
\label{ssec:mechanism}

Why is $s_7$'s partner integral when the dyadic floor of
Proposition~\ref{prop:dyadic} permits denominators? The answer is not that
$\eta$-quotients have integer Fourier coefficients. It is a normalization
statement, and the distinction is not pedantic --- the object in question is a
series \emph{reversion}, and reversion does not preserve integrality in
general.

Following~\cite{OBrien2016}, set
\begin{equation}\label{eq:obrien-objects}
  z_7 = \sum_{m,n\in\Z}q^{\,m^2+mn+2n^2}, \qquad
  Z_7 = \frac{7P(q^7)-P(q)}{6}, \qquad
  X_7 = \frac{\eta_1^3\eta_7^3}{z_7^{\,3}},
\end{equation}
with $P$ the Ramanujan Eisenstein series $P(q) = 1-24\sum_{n\geq1}\sigma(n)q^n$
and $\eta_1^3\eta_7^3 = q\prod_{j\geq1}(1-q^j)^3(1-q^{7j})^3$. O'Brien's
Theorem~6.1 gives the expansion $z_7 = \sum_{n\geq0}c_7(n)X_7^{\,n}$, and the
$c_7(n)$ are the terms of~\eqref{eq:s7-partner-recurrence}.

\begin{proposition}[Reversion of a normalized integral series]
\label{prop:normalized-reversion}
Let $X \in q\,\Z[[q]]$ have leading coefficient exactly $1$, i.e.
$X = q + O(q^2)$, and let $Y \in \Z[[q]]$. Then the unique expansion
$Y = \sum_{n\geq0}c_nX^n$ has $c_n \in \Z$ for all $n$.
\end{proposition}

\begin{proof}
Determine the $c_n$ recursively: $c_0 = Y(0)$, and having fixed
$c_0,\dots,c_{n-1}$, the residual $Y - \sum_{k<n}c_kX^k$ lies in
$q^n\Z[[q]]$, while $X^n = q^n + O(q^{n+1})$; so $c_n$ is the coefficient of
$q^n$ in the residual divided by $1$, hence an integer. Integrality of the
leading coefficient of $X$ is exactly what makes that division trivial; a
leading coefficient of $2$, say, would reintroduce denominators of precisely
the kind Proposition~\ref{prop:dyadic} permits.
\end{proof}

\begin{proposition}[The mechanism, verified; status (E)]\label{prop:mechanism}
With $X_7$ as in~\eqref{eq:obrien-objects}, exact rational arithmetic on
$q$-expansions gives
\[
  X_7 \;=\; q - 9q^2 + 30q^3 - 15q^4 + \cdots \in q\,\Z[[q]],
\]
so $X_7$ is a normalized integral uniformizer; and expanding $z_7$ in powers
of $X_7$ reproduces $1,2,22,336,6006,117348,2428272$, i.e. A279619.
Consequently, by Proposition~\ref{prop:normalized-reversion}, integrality of
$c_7$ follows from~\eqref{eq:obrien-objects} together with O'Brien's
Theorem~6.1 --- but from those taken \emph{jointly}, and not from the
integrality of an $\eta$-quotient alone.
\end{proposition}

\begin{remark}[Resolving the source's equation (3.15); status (E)]
\label{rem:315}
The pinned copy of~\cite{OBrien2016} renders the denominator of its equation
(3.15) ambiguously between $z_7$ and $Z_7$ and drops the exponent, so which
object is meant is not determined by the text. Three readings were tested in
exact arithmetic:
\[
  \frac{\eta_1^3\eta_7^3}{z_7^{\,3}}, \qquad
  \frac{\eta_1^3\eta_7^3}{Z_7^{\,3}}, \qquad
  \frac{\eta_1^3\eta_7^3}{Z_7^{\,2}} .
\]
\emph{All three} are normalized integral uniformizers, and all three therefore
produce integral expansion coefficients --- which is exactly why integrality
alone cannot resolve the ambiguity. They produce, respectively,
\[
  1,2,22,336,6006,117348,2428272; \qquad
  1,2,34,828,23662,739984,24521448; \qquad
  1,2,26,476,10206,239280,5942232,
\]
so only the first reproduces A279619. It is also the only weight-consistent
reading: $\eta_1^3\eta_7^3$ has weight $3$ and $z_7$ has weight $1$, so
$z_7^{\,3}$ is what makes the quotient weight $0$, as a uniformizer must be.
The computation additionally verifies $z_7^{\,2} = Z_7$ exactly, confirming
that the two symbols denote genuinely different objects. The resolution is
recorded so that the ambiguity need not be re-litigated.
\end{remark}

\begin{remark}[Scope]\label{rem:integrality-scope}
Proposition~\ref{prop:mechanism} establishes nothing about $s_{10}$ or
$s_{18}$. Their partners are non-integral by Theorem~\ref{thm:non-integral},
and no modular parametrization by a normalized integral uniformizer is known
for them --- but absence of a known parametrization is not a proof that none
exists, and no such negative claim is made here. In particular no level is
assigned to $s_{18}$ anywhere in this paper.
\end{remark}

\subsection{The modular coordinate}

The parameter of the $s_7$ family is a modular function for the same group
that Remark~\ref{rem:signature} identified from the exponent data.

\begin{proposition}[status (E), with a cited group-theoretic input]
\label{prop:hauptmodul}
Let $t$ be the generating function of OEIS A279618~\cite{OEIS} and let
$u = q^{-1}\bigl(E(q)/E(q^7)\bigr)^4$, with $E(q)=\prod_{n\geq1}(1-q^n)$, be
the classical $\Gamma_0(7)$ Hauptmodul; put $v = 1/u \in q\,\Z[[q]]$. Then,
in exact rational arithmetic on $q$-expansions:
\begin{enumerate}
\item $t$ is \emph{not} a M\"obius function of $v$ (the linear system solved
  from the first 6 orders admits no solution verifying on later orders);
\item $t$ \emph{is} of the form $(\alpha v^2+\beta v+\gamma)/(\delta
  v^2+\varepsilon v+\zeta)$, with coefficients solved linearly from the first
  12 orders and then verified on held-out orders 13 and 29 that were not used
  in solving (29 orders available in total);
\item rewriting $t$ as a degree-2 rational function of $u$ and solving for the
  involution constant $\kappa$ with $R(\kappa/u)\equiv R(u)$ yields
  $\kappa = 49$.
\end{enumerate}
Items (i) and (ii) give $[\Q(u):\Q(t)] = 2$, so $\Q(t)$ is the fixed field of
an element normalizing $\Gamma_0(7)$. The normalizer of $\Gamma_0(N)$ in
$\mathrm{PSL}_2(\mathbb{R})$ for prime $N$ is the Fricke extension
$\Gamma_0(N)+$ \cite{AtkinLehner1970} --- this step is cited, not re-derived
--- so that element is the Fricke involution and $t$ is a $\Gamma_0(7)+$
Hauptmodul. Item (iii) exhibits the Fricke constant $\kappa = 7^2$ directly.
\end{proposition}

The constant $\kappa$ is computed from the fitted coefficients, never assumed;
the checker's verdict requires it to come out $49$ but does not supply that
value to the fit. Three negative controls accompany the computation: feeding
it $v$ itself (a genuine $\Gamma_0(7)$ Hauptmodul) must return the verdict
``not plus'', and does; corrupting one coefficient of A279618 must break the
degree-2 verification, and does; and testing against the wrong level,
$v_5 = q\bigl(E(q^5)/E(q)\bigr)^6$ for $\Gamma_0(5)$, must fail every fit, and
does.

\begin{remark}[Three independent indicators of the level]
\label{rem:level-indicators}
The number $7$ has now appeared from three computations that share no code:
the implied Fuchsian signature of the $s_7$ operators, matching
$\Gamma_0(7)+$ (Remark~\ref{rem:signature}, status (E)); the Hauptmodul and
Fricke constant $\kappa = 49$ of Proposition~\ref{prop:hauptmodul} (status
(E)); and the $\sqrt7$ appearing uninvited in the numerically continued
monodromy, together with $\det G = -14$ (Numerical Claim~\ref{nc:monodromy}
and Proposition~\ref{prop:lattice}, status (N)/(E)). What that convergence
does and does not license is the subject of \S\ref{sec:framework}.
\end{remark}

% ----------------------------------------------------------------------------
% Generated-by: Opus 5 (Stream 1 paper agent) | Verified-by: recurrence and
% values vs Agora/Sequences/PartnerIntegrality.lean (partnerPair, s7/s10/s18_
% partner_values, partnerSeq_s7_recurrence, s7_partner_integral_pass7) and
% Agora/Axioms/OBrien2016.lean; dyadic baseline vs Agora/Sequences/FormalSqrt.
% lean (sqrtSeq_dyadic, dyadic_baseline_control); X7/readings/c7 vs Stream 2
% data/certificates/S7_PARTNER_MODULAR.json; Hauptmodul orders, held-out orders,
% kappa and controls vs HAUPTMODUL_S7_GAMMA07PLUS.json and
% checkers/check_s7_hauptmodul_gamma07plus.py. Claims #7-#10 of paper/PLAN.md §3.
% | Reviewed-by: pending T0 (Xavier)
% ----------------------------------------------------------------------------
